\paragraph{Conjugate-coordinate class times a fuzzy factor: the tadpole of the static scalar channel at \(q=0\), the profile of \(\sqrt{|g|}\) and the response of \(\vartheta\) (conditional; bounded obstruction).}
This paragraph supplements the two preceding ones, in which the terms that cancel the tadpole of \(U_4\) at \(q=0\), the second-order parts of \(U_2\) and of the first term, and the response of the fibre levels to \(\bar h\) were not computed. Keep the setting and notation of the paragraph on the response of the scale of \(K\), with the operators at \(c_3=\mu=0\), \(V=T\), equal weights, \(\varrho=1\), \(q\cdot y=q\cdot v=0\) and large \(\tau\), and the parametrisation of its forms: a term \(c\,t^2\) of the density enters the reduced form as \(\tfrac12ct^2\) after the phase average, and \(\sqrt{|g|}=\sqrt{|\bar g|}\,(1+\sigma_1+\sigma_1^2)\) to second order, with \(\sigma_1=-t/2m^2\) and \(\sqrt{|\bar g|}=m^{-4}\), the form that gives the term \(U_4\sigma_1^2\) of the preceding paragraph. In this channel \(\bar h^{11}\) is proportional to \(1-g\simeq2/\sinh^2\tau\), so the wave is along \(y^\mu y^\nu\) up to terms of that order; it is not the pure-trace direction \(\bar h^{\mu\nu}\propto\bar\gamma^{\mu\nu}\).

\emph{Profile and stationarity.} At \(q=0\) write the part of the density without derivatives of the wave, at fixed \(\chi\), as \(\sum_i(T_i\sigma_1+C_i\sigma_1^2)\), summed over its terms. A term \(w_i\sqrt{|g|}\) whose coefficient \(w_i\) does not respond to the wave has \(T_i=C_i=w_i\sqrt{|\bar g|}\). This holds exactly for \(U_4\) and for any term of the cosmological type; for \(U_2\) and the fibre-constant part of the first term, whose derivative in \(\ln\chi\) is the fibre-constant part of \(Y\), the recorded rigidity gives \(T_i=w_i\sqrt{|\bar g|}\) at first order, and the tree potential \(\rho_MV_{\mathrm t}/g^2\) carries \(\rho_M\), which does not respond at fixed \(\chi\). The background is stationary in this channel at \(q=0\) exactly when \(\sum_iT_i=0\). With the exchange of \(\chi\) of the preceding paragraph, for \(q^2\) below \(m_\chi^2\), the denominator of the first term of the vacuum-term display is \(q^2+(L/2m^4)\sum_iC_i+q_*^2|_\chi\), the term \(C_i=U_4\) giving \(-q_b^2\); hence at a stationary background
\[
q^2-q_b^2+q_*^2=q^2+\frac{L}{2m^4}\sum_i\bigl(C_i-T_i\bigr)+q_*^2\big|_\chi ,
\]
in which every term carrying \(\sqrt{|g|}\) alone drops out. A redefinition \(\sigma_1\to\sigma_1+a\sigma_1^2\) shifts each \(C_i\) by \(aT_i\), so the quadratic part of a term with an uncancelled tadpole depends on the parametrisation, while the right-hand side at \(\sum_iT_i=0\) does not. So no set of terms carrying \(\sqrt{|g|}\) alone, a term of the cosmological type included, cancels the tadpole at \(q=0\) and leaves a non-derivative part: where \(U_4\) and such a term are the only ones, the cancellation removes \(-q_b^2\) with the tadpole. A term of that type that does not depend on \(\chi\) leaves \(U'_{\mathrm{loop}}\), the stationarity in \(\chi\), \(\beta_K\) and \(q_*^2|_\chi\) unchanged, so the requirement recorded above, that the vacuum term be cancelled at the minimum without cancelling its derivative in \(\chi\), is compatible with this identity. Since \(q_*^2|_\chi=-L\,U_{\mathrm{loop}}'^2/(4m^4U_{\chi\chi})<0\) at a local minimum with \(U'_{\mathrm{loop}}\ne0\), a Yukawa pole of this term at a stationary background needs \((L/2m^4)\sum_i(C_i-T_i)>-q_*^2|_\chi\) from terms that do not carry \(\sqrt{|g|}\) alone: within the recorded content these are the second-order parts of \(U_2\) and of the fibre-constant part of the first term, the part of the first term with the fibre flux, at first and at second order, and the tree non-derivative terms of the wave, which the leading symbol of the tree form drops, none of which is computed. With them set to zero the pole of this term has the tachyonic sign for either sign of \(\hat U\), so at a stationary background the sign of \(\hat U\) does not decide it.

\emph{The response of \(\vartheta\).} In the channel above \(\vartheta\) does not respond at first order: it enters \(U_2\) and the fibre-constant part of the first term through the Gaussian average over the loop momentum, which a perturbation along \(y^\mu y^\nu\) leaves unchanged since \(\bar T^{ab\mu}y_\mu=0\). Reading \(\vartheta=-y\cdot y/m^2\) in the fluctuated metric at fixed \(y\) and \(m\) would move it there, by \(\delta\ln\vartheta=-\epsilon\,y\cdot y\) for \(\delta\gamma^{\mu\nu}=\epsilon\,y^\mu y^\nu\), so that reading is not the response of these densities. Along the pure-trace direction the rigidity does not apply and the response is not computed. Take as response hypothesis that there \(\delta\ln\vartheta=p\,\sigma_1\) at first order, with \(p\) a constant and \(\mu_J\), \(c_J\), \(M_A^2\) held; \(p=0\) extends the rigidity, and the reading just excluded in the channel above would give \(p=\tfrac12\). Integrating the level form of \(Y\) in \(\ln\chi\) with the proper-time integral that produces it, with a lower limit \(s_\Lambda\) as in the preceding paragraph, the first term is the density
\[
-\frac{\sqrt{|g|}}{32\pi^2}\sum_A\sum_{J\le n}(2J+1)\Bigl[\tfrac{15}{4}\vartheta^2E_1(s_\Lambda\mathcal D)-\frac{\vartheta c_J}{\mathcal D}+\frac{c_J^2}{\mathcal D^2}\Bigr],
\]
with \(E_1\) the exponential integral of order one and up to terms of relative order \(s_\Lambda\mathcal D\); the logarithm of \(E_1\) at small \(s_\Lambda\mathcal D\) drops out of \(Y\). Under the hypothesis the first-order part of this density along that direction exceeds the profile of \(\sqrt{|g|}\) by \(\sigma_1\) times \(-p\,(\sqrt{|g|}/32\pi^2)\sum_A\sum_{J\le n}(2J+1)[\tfrac{15}{2}\vartheta^2E_1(s_\Lambda\mathcal D)-\vartheta c_J/\mathcal D]\), whose value depends on \(s_\Lambda\) through \(E_1\), the bracket at each level being positive exactly when \(E_1(s_\Lambda\mathcal D)>2c_J/(15\vartheta\mathcal D)\), and that of \(U_2\) exceeds it by \(p\,U_2\sigma_1\); the coefficient of \(\sigma_1\) in the response of \(Y\) is finite, with the summand \((\sqrt{|g|}/32\pi^2)(2J+1)\,2M_A^2[(1+2p)\tfrac{15\vartheta^2}{4\mathcal D}-(1+p)\tfrac{\vartheta c_J}{\mathcal D^2}+\tfrac{2c_J^2}{\mathcal D^3}]\). These parts do not enter the identity above, which is for the channel along \(y^\mu y^\nu\).

These statements are for the hypotheses of the paragraph on the response of the scale of \(K\), with the operators at \(c_3=\mu=0\) and \(V=T\), at \(q=0\) in its static channel at large \(\tau\) and in the parametrisation stated, and, along the pure-trace direction, under the response hypothesis stated, with the levels held and at the leading order in \(s\) of the level form. The identity at a stationary background is exact for these forms, and the density of the first term is exact for the level form at the stated limit. They do not compute the second-order parts of \(U_2\) and of the first term, the response of the fibre levels or of the flux insertions to \(\bar h\), the term of \(U_2\) proportional to \(W_R\), the tree non-derivative terms, higher loops or the fibre Laplacian, do not determine \(\mathcal M\), \(p\) or the sign of the static scalar pole, and do not identify \(t\) with a quantity of the emergent metric of Chapter~\ref{ch:vierbein-emergence}, so no Cassini or Newtonian comparison follows. Problems~C, F and~N remain open and PROGRAMME-CONDITIONAL.